\chapter*{不定积分每日一题}

\refer{《积分的方法与技巧》金玉明\ 顾新身\ 毛瑞庭\ /\ 编著}

\exercise{求积分$\int{\cot^n{x}\d x}$}%在花括号内输入题目

\teacher{张晓鹏}%在花括号内输入你的名字

\bigresult{
    \formula{\int{\cot^{2m+1}{x}\d x}=\sum_{r=0}^{m-1}\qty(-1)^{r+1}\frac{\cot^{2m-2r}{x}}{2m-2r}+\qty(-1)^m\ln{\abs{\csc{x}}}+C}
    \formula{\int{\cot^{2m}{x}\d x}=\sum_{r=0}^{m-1}\qty(-1)^{r+1}\frac{\cot^{2m-2r-1}{x}}{2m-2r-1}+\qty(-1)^mx+C}
}%在花括号内输入本题结果

\analysis{
    $\cot^n{x}$的积分和$\tan^n{x}$方法差不多，也是主要靠下面这个递推公式.
    \formula{
        \begin{aligned}
            \int{\cot^{x}\d x}&=\int{\cot^{n-2}{x}\qty(\csc^2{x}-1)\d x}\\
            &=-\int{\cot^{n-2}{x}\d\cot{x}}-\int{\cot^{n-2}}{x}\d x\\
            &=\frac{\cot^{n-1}x}{n-1}-\int{\cot^{n-2}{x}\d x}
        \end{aligned}
    }
    加上起始项
    \formula{\int{\cot{x}\d x}=\int{\frac{\d\sin{x}}{\sin{x}}}=\ln{\abs{\sin{x}}}+C}
    \formula{\int{\cot^2{x}\d x=\int{\qty(\csc^2{x}-1)\d x}=-\cot{x}-x+C}}
    即可递推得到最后结果.
    \formula{\int{\cot^{2m+1}{x}\d x}=\sum_{r=0}^{m-1}\qty(-1)^{r+1}\frac{\cot^{2m-2r}{x}}{2m-2r}+\qty(-1)^m\ln{\abs{\csc{x}}}+C}
    \formula{\int{\cot^{2m}{x}\d x}=\sum_{r=0}^{m-1}\qty(-1)^{r+1}\frac{\cot^{2m-2r-1}{x}}{2m-2r-1}+\qty(-1)^mx+C}
}%在花括号内输入分析